% Generated draft: requires mathematical and source review.
\documentclass[11pt]{article}
\usepackage[margin=25mm]{geometry}
\usepackage{fontspec}
\setmainfont{Latin Modern Roman}
\usepackage{amsmath,amssymb,mathtools,unicode-math}
\setmathfont{Latin Modern Math}
\usepackage{enumitem,needspace,xurl,hyperref}
\setlength{\parindent}{0pt}
\setlength{\parskip}{5pt}
\setlength{\emergencystretch}{3em}
\newcommand{\N}{\mathbb{N}}
\newcommand{\Z}{\mathbb{Z}}
\newcommand{\Q}{\mathbb{Q}}
\newcommand{\R}{\mathbb{R}}
\newcommand{\C}{\mathbb{C}}
\newcommand{\F}{\mathbb{F}}
\newcommand{\A}{\mathbb{A}}
\newcommand{\PP}{\mathbb{P}}
\newcommand{\E}{\mathbb{E}}
\newcommand{\eps}{\varepsilon}
\newcommand{\dd}{\,\mathrm{d}}
\begin{document}
\section*{UnsolvedMath ID: 10400125}
\textbf{Source problem code:} catalogue entry
\subsection{Definitions and mathematical statement}
For the source's quantum invariant \(Z_r^G(M)\), let \(d_0,\ldots,d_n\) be its indicated finite list of exponents and put \(d=\max_{0\leq j\leq n}d_j\). The growth rate conjecture is
\[
 |Z_r^G(M)|=O(r^d)\qquad(r\to\infty).
\]
Explicitly, for each fixed admissible \(M,G\), there should be constants \(C<\infty\) and \(r_0\) such that
\[
 |Z_r^G(M)|\leq Cr^d\qquad\text{for all admissible levels }r\geq r_0.
\]
This is an upper bound; it does not assert a matching lower bound, a limiting leading constant, or absence of oscillatory cancellations.

The present excerpt does not define the \(d_j\). The adjacent local Conjecture 7.7 identifies them as rational exponents in a proposed asymptotic sum over the finitely many Chern--Simons values of flat connections on a closed oriented 3-manifold, with sector terms proportional to \(e^{2\pi i rc_j}r^{d_j}\). That expansion is itself conjectural and is not assumed proved in order to state the separate growth bound. A further topological formula is a separate Conjecture 7.9, not an established definition substituted here.

The source additionally reports that some polynomial-in-level upper bound is known, with exponent bounded by a simple \(G\)-dependent function of Heegaard genus. Heegaard genus is the least genus of a surface splitting the closed 3-manifold into two handlebodies. No explicit genus function is supplied. This coarse known-growth report does not identify the specific exponent \(d\) or settle the conjectured bound.

Normalization and level conventions remain to be recovered: the adjacent expansion uses \(Z_k^G\) and a shifted variable \(r=k+h^\vee\), whereas this excerpt prints \(Z_r^G\). The admissible \(G\), invariant and parameter definitions are missing.
\subsection{Short English statement}
Bound the quantum invariant by the maximal power predicted in its Chern--Simons asymptotic expansion. A general polynomial-growth bound is reported, but the exact exponents, normalization, and level conventions need their source definitions.
\subsection{Sources}
\begin{itemize}
\item Local UnsolvedMath registry, permanent ID 10400125, source code catalogue entry; no proposer or proposal year supplied; updated 2026-08-06.
\item Ohtsuki, \emph{Problems on invariants of knots and 3-manifolds} (2002), Conjecture 7.10, PDF page 107: \url{https://msp.org/gtm/2002/04/gtm-2002-04-024s.pdf}. Registry access date 2026-07-29; PDF font-structure extraction. Provenance distinguishes 231 numbered bold headings. Present status and redistribution terms require review.
\item Local context: ID 10400122, source code AMR-103-0122, Conjecture 7.7, PDF page 106, supplies the proposed asymptotic-exponent context and credits J. E. Andersen, reference [6]. ID 10400124, Conjecture 7.9, separately proposes a topological interpretation. Full bibliography entries are absent; no external source was visited.
\item The supplied summary repeats the Heegaard-genus polynomial bound without a precise formula or resolving citation.
\end{itemize}
\end{document}
